\documentclass[11pt,a4paper]{article}

\usepackage{fontspec}
\usepackage{xeCJK}
\setCJKmainfont{Hiragino Mincho ProN}
\setCJKsansfont{Hiragino Kaku Gothic ProN}
\usepackage[margin=2.5cm]{geometry}
\usepackage{amsmath,amssymb}
\usepackage{hyperref}
\usepackage{booktabs}
\usepackage{amsthm}
\usepackage{tcolorbox}

\tcbuselibrary{breakable,skins}

\newtcolorbox{epigraph}[1][]{
  colback=white, colframe=black!30,
  fonttitle=\itshape, title={},
  breakable, sharp corners, boxrule=0.5pt,
  left=15pt, right=15pt
}

\newtheorem{theorem}{Theorem}
\newtheorem{observation}[theorem]{Observation}
\newtheorem{conjecture}[theorem]{Conjecture}

\newcommand{\Spec}{\mathrm{Spec}}
\newcommand{\Sh}{\mathrm{Sh}}
\newcommand{\Zp}{\mathbb{Z}[1/p]}

\setlength{\parskip}{0.6em}
\setlength{\parindent}{0pt}

\begin{document}

\begin{center}
\vspace*{1cm}
{\Huge\bfseries 点の背後で蠢くもの}\\[0.8cm]
{\Large\bfseries グロタンディークの層の圏で時空を書き換える}\\[0.5cm]
{\large 算術的真空工学の理論・実験・展望}\\[2cm]
{\large Wright Brothers}\\[0.5cm]
{\normalsize 独立研究}\\[0.5cm]
{2026年4月}\\[3cm]
{\itshape アレクサンドル・グロタンディーク（1928--2014）の記憶に捧げる}
\vspace*{1cm}
\end{center}

\newpage
\tableofcontents
\newpage

% ============================================================================
%  EPIGRAPH
% ============================================================================

\begin{epigraph}
\textit{La chose la plus importante dans un mathématicien, ce n'est pas le problème
qu'il résout, mais la vision qu'il ouvre.}

（数学者にとって最も重要なのは、彼が解いた問題ではなく、彼が開いた視界である。）

\hfill --- Alexander Grothendieck
\end{epigraph}

\vspace{1cm}

% ============================================================================
\part{パラダイムの転換}
% ============================================================================

\section{カントールからグロタンディークへ}

\subsection{点の実体視：カントールの遺産}

19世紀末、ゲオルク・カントールは集合論を創始し、
数学の基礎を「点の集合」の上に置いた。
空間は点の集まりであり、関数は点の上の値の対応であり、
物理量は点の上で定義された数値である。

20世紀の物理学はこの基盤の上に建てられた。
量子力学のヒルベルト空間は「状態ベクトルの集合」、
場の量子論は「時空の各点での場の値の集合」、
一般相対論は「多様体上の点の集合に計量を載せたもの」。

真空エネルギーも例外ではない:
$$E = \frac{\hbar\omega_0}{2}\sum_{n=1}^{\infty} n$$

ここで各 $n$ は「点」（モード番号）であり、
和は「点の集合」上の足し上げである。
カントール的視点: 点を足す。それだけ。

\subsection{グロタンディークの革命：点の背後}

1960年代、アレクサンドル・グロタンディークは
数学を根底から覆した。

\begin{epigraph}
\textit{空間の本質は「点の集合」ではない。
空間の本質は、その上に載っている「層の圏」にある。}
\end{epigraph}

グロタンディークは問うた: 点とは何か？
答え: 点は\textbf{射}（morphism）である。
$\Spec(\mathbb{Z})$の点$(p)$は「スキームから$(p)$への射」であり、
点自体が独立した実体ではなく、\textbf{構造全体との関係}で定義される。

これは認識論的転換である。
カントール: 点が先にあり、空間はその集まり。
グロタンディーク: \textbf{構造}が先にあり、点はその影。

\subsection{物理学への翻訳}

この転換を物理に持ち込む:

\textbf{カントール的物理学}（現在の標準）:
$$\zeta(s) = \sum_{n=1}^{\infty} n^{-s} \quad \text{（点の足し上げ）}$$

モードは「点」。真空エネルギーは「点の和」。
モードを消す = 和から項を引く = 量的変化。

\textbf{グロタンディーク的物理学}（本研究）:
$$\zeta(s) = \text{$\Spec(\mathbb{Z})$のゼータ関数} \quad \text{（空間の不変量）}$$

$\zeta(s)$は$\Spec(\mathbb{Z})$という\textbf{空間自体の性質}。
「モードを消す」= 空間を$\Spec(\Zp)$に変える = \textbf{質的変化}。

和の項を引くのではなく、\textbf{空間そのものを取り替える}。
鍵を壊すのではなく、\textbf{鍵穴ごと消す}。

\section{直観主義論理とGaussian Splatting}

\subsection{排中律が崩れる世界}

古典論理: 命題は真か偽かのどちらか（排中律 $P \lor \neg P$）。
物理法則は「真」か「偽」。WECは「成立する」か「しない」か。

グロタンディークのトポス理論: 真偽値は$\{0, 1\}$ではなく
\textbf{ハイティング代数}（直観主義論理の真偽値の空間）。
排中律は一般には成り立たない。
命題は「真」「偽」の他に「部分的に真」「文脈に依存」があり得る。

Döring-Isham~\cite{DoeringIsham2008}はこれを量子力学に適用:
量子系の命題は、観測の文脈（同時に測定可能な物理量の組）に依存して
異なる真偽値を持つ。コッヘン-シュペッカーの定理がこれを保証する。

\subsection{Gaussian Splattingのメタファー}

\begin{epigraph}
3Dコンピュータグラフィクスの最新技術「Gaussian Splatting」は、
シーンを「ポリゴン（硬い面の集合）」ではなく
「ガウシアンの重ね合わせ（柔らかい雲の集合）」で表現する。

同じ3D空間でも、\textbf{「何でレンダリングするか」}によって
見え方が全く異なる。

これはグロタンディークの層の圏と同じ発想:
\textbf{同じ「空間」でも、載せる「層」（= レンダリング方法）を変えると、
まったく異なる「世界」が現れる。}
\end{epigraph}

物理への翻訳:
\begin{itemize}
\item $\Sh(\Spec(\mathbb{Z})_{\text{ét}})$でレンダリング → WEC = 真（通常の宇宙）
\item $\Sh(\Spec(\Zp)_{\text{ét}})$でレンダリング → WEC = 偽（ワープ可能な宇宙）
\end{itemize}

「宇宙」は同じ。変わるのは「レンダリングエンジン」（トポス）。
局所化$\mathbb{Z} \to \Zp$はレンダリングエンジンの切り替え。

\textbf{物理法則の書き換えが可能なのは、
物理法則が「絶対的真理」ではなく
「レンダリング結果」だからである。}

% ============================================================================
\part{実験：ζ(-1)からζ(-3)へ}
% ============================================================================

\section{第1段階：ζ(-1)の符号反転}

\subsection{実験の核心}

1Dのゼータ正則化真空エネルギー:
$E \propto \zeta(-1) = -1/12$。

偶数モードを抑制: $E \propto \zeta_{\neg 2}(-1) = +1/12$。

符号が反転する。
これを超伝導キャビティ + SQUIDノッチフィルタで測定する。

\subsection{装置（48万円構成）}

\begin{itemize}
\item 3Dアルミキャビティ（35 × 35 × 10 mm）: 基本モード 6.06 GHz
\item Al/AlO$_x$/Al SQUID接合 1個: 第2高調波に同調
\item $^3$He冷凍機: 300 mK
\item VNA: $S_{21}$測定
\end{itemize}

決定的測定: SQUID磁束スイープ中のAC Starkシフトの挙動
→ 不連続ジャンプ（レギュレータ予想支持）or 滑らかな変化（反証）。

\section{第2段階：ζ(-3)の符号反転 = ワープバブルの検証}

\subsection{1Dから3Dへ}

$\zeta(-1)$は1D真空エネルギー。ワープに必要なのは3D:
$$\zeta(-3) = \frac{1}{120} \quad \to \quad \zeta_{\neg 2}(-3) = \frac{1-8}{120} = -\frac{7}{120}$$

$\zeta(-1)$の符号反転が確認されたら、次は$\zeta(-3)$。

\subsection{3Dキャビティへの拡張}

1Dの伝送線共振器 → 3Dのマイクロ波キャビティ。

3D矩形キャビティのゼロ点エネルギー:
$$E_{3D} = \frac{\hbar}{2}\sum_{m,n,p} \omega_{mnp}
= \frac{\hbar c}{2}\sum_{m,n,p} \pi\sqrt{(m/a)^2 + (n/b)^2 + (p/d)^2}$$

この3重和のゼータ正則化が$\zeta(-3)$に対応する
（エプシュタイン・ゼータ関数を経由）。

\subsection{必要な追加装置}

\begin{itemize}
\item 3D長方形キャビティ（異方的: $a \neq b \neq d$）
  → モードの縮退を解いて個別に操作可能にする
\item SQUID 3--5個: 偶数モード$(2,0,1)$, $(0,2,1)$, $(2,2,1)$, ...を各々抑制
\item 精密な磁束制御: 各SQUIDに独立したバイアス線
\end{itemize}

コスト追加: ~50万円（SQUID追加 + キャビティ再設計）。
合計: ~100万円。

\subsection{成功の意味}

$\zeta_{\neg 2}(-3) < 0$が実験で確認されたら:
\begin{itemize}
\item 3Dの真空エネルギー密度が負 = WEC違反の物理的実現
\item アルクビエレ計量の「材料」が存在することの直接証拠
\item → ワープドライブの原理的可能性が実験的に確立される
\end{itemize}

% ============================================================================
\part{驚きの一致：偶然か構造か}
% ============================================================================

\section{48と251：異種の不変量が語ること}

この研究プログラムの中で、最も「こじつけ」から遠い証拠は
全世代レプトンの$g-2$異常の算術的一致である。

\begin{observation}[電子: K理論]
$$\Delta a_e = |K_3(\mathbb{Z})| \times 10^{-14} = 48 \times 10^{-14}$$
実験値: $\sim 4.8 \times 10^{-13}$。一致。
\end{observation}

\begin{observation}[ミューオン: ゼータ値]
$$\Delta a_\mu = \left(\frac{1}{|\zeta(-5)|} - 1\right) \times 10^{-11} = 251 \times 10^{-11}$$
実験値: $(251 \pm 59) \times 10^{-11}$。中心値完全一致。
\end{observation}

48はK理論の不変量。251はゼータ関数の特殊値。
これらは$\Spec(\mathbb{Z})$の\textbf{代数的に独立な}不変量である。

\begin{epigraph}
2つの無関係な数学的構造（K群とゼータ値）が、
2つの無関係な物理量（電子とミューオンの磁気モーメント）を
同時に、チューニング不能な整数として再現する。

これを偶然と呼ぶには、宇宙があまりに精密である。
\end{epigraph}

\section{137 = 120 + 12 + 5}

微細構造定数$\alpha^{-1} = 137.036$の分解:
$$\alpha^{-1} = \underbrace{\frac{1}{\zeta(-3)}}_{120} + \underbrace{\frac{1}{|\zeta(-1)|}}_{12} + \underbrace{g(132)}_{5} + 4(\zeta(6) - \zeta(7))$$

\begin{itemize}
\item 120 = 3Dカシミールエネルギーの分母
\item 12 = 1Dカシミールエネルギーの分母 = 標準模型ゲージ群の次元
\item 5 = 合成数132から次の素数137への素数ギャップ
\item $4(\zeta(6)-\zeta(7))$ = 時空次元$d=4$での高次元補正
\end{itemize}

自由パラメータ: ゼロ。精度: 0.00002\%。

\section{π(137) = 33: 自己参照する宇宙}

137以下の素数の個数は33。
この数は3箇所に独立に出現する:
\begin{enumerate}
\item $\ln(\Lambda_{\mathrm{GUT}}/m_Z) \approx 33$（繰り込み群の走り）
\item $\Delta m^2_{32}/\Delta m^2_{21} \approx 33$（ニュートリノ質量比）
\item $\pi(137) = 33$（素数計数関数の定義）
\end{enumerate}

$\alpha$の値が、$\alpha$自身以下の素数の個数を参照して決まる。
宇宙は\textbf{自己参照的}。

\begin{epigraph}
ゲーデルの不完全性定理が「数学は自分自身について語る時に限界に達する」
ことを示したように、$\alpha = f(\pi(\alpha^{-1}))$は
「宇宙は自分自身の構造を参照して物理定数を決めている」ことを示唆する。

$\Spec(\mathbb{Z})$は自己参照するコード。
物理定数は、コードが自分自身をコンパイルした結果。
\end{epigraph}

% ============================================================================
\part{展望：何が可能になるか}
% ============================================================================

\section{$\zeta(-1)$の先にある世界}

第1段階（$\zeta(-1)$符号反転）が成功したとする。
そこから開ける道:

\textbf{ワープドライブ}（$\zeta(-3) < 0$の検証後）:
WEC違反を空間的に球殻状に配置 → アルクビエレ計量。
エネルギーコスト: レギュレータ予想が正しければ$\sim 10^{12}$ J。

\textbf{ワームホール}（$\zeta(-3) < 0$ + トロイダル幾何）:
2つのワープバブルの壁をトンネルで接続。
算術的ドメインウォール$\Spec(\mathbb{F}_p)$が喉の構造を提供。

\textbf{全世代時空工学}（多重素数コンソール）:
$N$個の素数チャンネルを独立制御 → $2^N$通りの真空状態。
各世代のレプトンの真空結合を選択的に操作。
$N = 20$で$\sim 10^6$状態の「時空プログラミング」。

\textbf{カルダシェフを超えて}:
エネルギーの量ではなく、物理法則の制御範囲で文明を測る。
Type A-0（法則を読む）→ A-1（変える）→ A-2（プログラムする）。

\section{タウの予測：未来への手紙}

$$\Delta a_\tau = (1/|\zeta(-9)| - 1) \times 10^{-8} = 131 \times 10^{-8}$$

131は素数。この予測は現在の実験精度では検証できない。
将来のスーパータウチャームファクトリーでのみ測定可能。

もしこの予測が確認されれば、3つのレプトン世代全てが
$\Spec(\mathbb{Z})$の異なる不変量（K群、$\zeta(-5)$、$\zeta(-9)$）
で統治されていることの完全な証拠になる。

\begin{epigraph}
未来の実験物理学者へ:
もしあなたがタウの異常磁気モーメントを測定し、
$131 \times 10^{-8}$という値を得たなら、
それは$\Spec(\mathbb{Z})$が時空の構造であることの
3番目の、そして最終的な確認である。

我々はそれを2026年に予測した。
\end{epigraph}

% ============================================================================
\part{グロタンディークへ}
% ============================================================================

\section{見えるものの背後にあるもの}

グロタンディークは2014年に亡くなった。
晩年はピレネー山中で隠遁生活を送り、
世俗から離れて思索を続けていた。

彼が生涯をかけて追い求めたのは、
数学の表層（個々の定理、具体的な計算）の背後にある
\textbf{「構造」}だった。

層（sheaf）、スキーム（scheme）、トポス（topos）。
これらの概念は全て、「見えるもの」の背後にある
「見えない構造」を言語化するためのものだった。

\begin{epigraph}
\textit{La question n'est pas de savoir
si l'univers est mathématique,
mais quelle mathématique il choisit.}

（問いは「宇宙が数学的か」ではなく、
「宇宙がどの数学を選んでいるか」である。）
\end{epigraph}

本研究は、グロタンディークの数学が
「美しいが抽象的な理論」ではなく
「宇宙の具体的な構造」であるという仮説を提示した。

$\Spec(\mathbb{Z})$は抽象代数の対象物ではなく、
真空が実際にその上に載っている空間かもしれない。
層のコホモロジーは数学者の道具ではなく、
物理法則が記述される言語かもしれない。
トポスの内部論理は論理学の遊びではなく、
WECの真偽が決まるメカニズムかもしれない。

48万円の実験がこの「かもしれない」を
「である」に変える可能性を持っている。

\section{最後に}

グロタンディークは、層の圏を「見えるものの背後で蠢くもの」
として導入した。

我々は、その「蠢くもの」の中に
物理法則の起源を見出し、
それを操作する方法を考案し、
48万円でテストする実験を設計した。

もし実験が成功すれば、
人類は初めて「物理法則を読む」存在から
「物理法則を書く」存在へと変貌する。

それは数学者グロタンディークが夢見た世界——
表層の背後にある構造が全てを支配する世界——
が物理的現実であることの確認になる。

\begin{epigraph}
\textit{Un mathématicien qui n'est pas un peu poète
ne sera jamais un vrai mathématicien.}

（少しも詩人でない数学者は、真の数学者にはなれない。）

\hfill --- Alexander Grothendieck
\end{epigraph}

\vfill

\begin{thebibliography}{15}
\bibitem{DoeringIsham2008} A.~Döring and C.~J.~Isham, J. Math. Phys. \textbf{49}, 053515 (2008).
\bibitem{Connes1994} A.~Connes, \emph{Noncommutative Geometry}, Academic Press (1994).
\bibitem{ConnesMarcolli2008} A.~Connes and M.~Marcolli, \emph{NCG, Quantum Fields and Motives}, AMS (2008).
\bibitem{BostConnes1995} J.-B.~Bost and A.~Connes, Selecta Math. \textbf{1}, 411 (1995).
\bibitem{Grothendieck} A.~Grothendieck, \emph{Récoltes et Semailles} (1985, unpublished autobiography).
\end{thebibliography}

\end{document}
